% Source: source/普通高考/2012/2012湖南理.pdf, page 1, question 14.
% Vector flowchart; pdfimages -list reports no embedded image for this page.
% Nodes: 开始; 输入 x,n; S=6; i=n-1; i>=0?; S=S·x+i+1; i=i-1; 输出 S; 结束.
% Directed edges: 开始→输入→S=6→i=n-1→判定;
% 判定 --是→ 累加→i=i-1→回到判定前的主干;
% 判定 --否→输出 S→结束. The loop return is separate from the output branch.
\begin{tikzpicture}[
  >=Stealth,
  flow/.style={draw, line width=0.55pt, -{Stealth[length=2.1pt,width=1.8pt]}},
  every node/.style={align=center, inner sep=0pt,
    execute at begin node=\setlength{\parindent}{0pt}},
  box/.style={draw, rectangle, minimum height=0.68cm,
    inner xsep=0.16cm, inner ysep=0.08cm, align=center,
    execute at begin node=\setlength{\parindent}{0pt}},
  terminal/.style={box, rounded corners=0.18cm, minimum width=1.40cm},
  decision/.style={draw, diamond, aspect=2.7, minimum width=3.05cm,
    minimum height=1.05cm, inner xsep=0.10cm, inner ysep=0.06cm,
    align=center, execute at begin node=\setlength{\parindent}{0pt}},
  io/.style={box, trapezium, trapezium left angle=72, trapezium right angle=108,
    minimum width=2.35cm},
]
  \path[use as bounding box] (-2.75,0) rectangle (5.85,12.55);

  \node[terminal] (start) at (0,12.05) {开始};
  \node[box, minimum width=3.20cm] (input) at (0,10.65) {输入 $x,n$};
  \node[box, minimum width=1.70cm] (sinit) at (0,9.25) {$S=6$};
  \node[box, minimum width=3.05cm] (iinit) at (0,7.95) {$i=n-1$};
  \node[decision] (test) at (0,4.15) {$i\geq 0?$};
  \node[box, minimum width=4.65cm] (update) at (3.10,5.35)
    {$S=S\cdot x+i+1$};
  \node[box, minimum width=3.00cm] (decrement) at (3.10,7.00) {$i=i-1$};
  \node[io, minimum width=2.35cm] (output) at (0,2.25) {输出 $S$};
  \node[terminal] (finish) at (0,0.75) {结束};
  \coordinate (loopjoin) at (0,7.35);

  \draw[flow] (start) -- (input);
  \draw[flow] (input) -- (sinit);
  \draw[flow] (sinit) -- (iinit);
  \draw[-] (iinit) -- (loopjoin);
  \draw[flow] (loopjoin) -- (test.north);
  \draw[flow] (test.east) -- node[above, inner sep=1pt] {是} (3.10,4.15)
    -- (update.south);
  \draw[flow] (update.north) -- (decrement.south);
  \draw[flow] (decrement.north) -- (3.10,7.35) -- (loopjoin);
  \draw[flow] (test.south) -- node[right, inner sep=1pt] {否} (output.north);
  \draw[flow] (output) -- (finish);
\end{tikzpicture}
